\documentclass[10pt,aspectratio=169]{beamer}
\usepackage{../../Common/course_slides}
\usepackage{../../Common/course_macros}
\usepackage[spanish,provide=*]{babel}

\title[Semana 1: Construcción del Arsenal Indicial]{De la Crisis Cartesiana al Arsenal de Índices}
\subtitle{Construyendo la notación de Einstein, $\delta_{ij}$ y $\epsilon_{ijk}$ desde la geometría}
\author{Geometría y Física I}
\institute{Licenciatura en Física}
\date{}

\begin{document}
\setlength{\abovedisplayskip}{3pt}
\setlength{\belowdisplayskip}{3pt}

% ---------------------------------------------------------------------
\begin{frame}
  \titlepage
\end{frame}

% ---------------------------------------------------------------------
\begin{frame}{Actividad 1: El Dolor Cartesiano y cómo Extraer Componentes}
  \begin{block}{Reto en clase (5 minutos)}
    \small
    Denotemos los vectores unitarios como $\hat{e}_1 = \hat{i}$, $\hat{e}_2 = \hat{j}$, $\hat{e}_3 = \hat{k}$. Todo vector se expande como $\vec{A} = \sum_{j=1}^3 A_j\hat{e}_j$.
    \begin{enumerate}
      \item[(a)] Si calculas a mano la primera componente de $\vec{A} \times (\vec{B} \times \vec{C})$, obtienes:
      \[
      [\vec{A} \times (\vec{B} \times \vec{C})]_1 = B_1(A_2 C_2 + A_3 C_3) - C_1(A_2 B_2 + A_3 B_3).
      \]
      ¿Qué término debes sumar y restar artificialmente para completar los productos punto $\vec{A}\cdot\vec{C}$ y $\vec{A}\cdot\vec{B}$?
      \item[(b)] Queremos extraer la componente $A_i$ de $\vec{A}$ proyectando con $\hat{e}_i$: $\hat{e}_i \cdot \vec{A} = \sum_{j=1}^3 A_j (\hat{e}_i \cdot \hat{e}_j)$.
      ¿Cuánto vale el producto $\hat{e}_i \cdot \hat{e}_j$ cuando $i = j$ y cuando $i \neq j$?
      \item[(c)] Usa tu respuesta de (b) para simplificar $\vec{A}\cdot\vec{B} = \left(\sum_i A_i\hat{e}_i\right)\cdot\left(\sum_j B_j\hat{e}_j\right)$.
    \end{enumerate}
  \end{block}
\end{frame}

% ---------------------------------------------------------------------
\begin{frame}{Discusión 1: Nacimiento de Einstein y de la Delta $\delta_{ij}$}
  \begin{columns}[c]
    \begin{column}{0.52\textwidth}
      \begin{block}{Construcción Geométrica de $\delta_{ij}$}
        No es una matriz arbitraria; es la \textbf{tabla de multiplicar de la base ortonormal}:
        \[
        \mathbf{\delta_{ij} \equiv \hat{e}_i \cdot \hat{e}_j} = \begin{cases} 1 & \text{si } i = j \\ 0 & \text{si } i \neq j \end{cases}
        \]
      \end{block}
      \small
      \textbf{Convenio de Einstein:} Todo índice repetido dos veces en un término se suma automáticamente de $1$ a $3$:
      \[
      A_i = \hat{e}_i \cdot (A_j\hat{e}_j) = A_j(\hat{e}_i\cdot\hat{e}_j) \implies \mathbf{\delta_{ij} A_j = A_i}
      \]
    \end{column}
    \begin{column}{0.45\textwidth}
      \begin{exampleblock}{El Interruptor Selector}
        \small
        \begin{itemize}
          \item Contraer con $\delta_{ij}$ apaga los términos con $j \neq i$ y sustituye $j \to i$.
          \item \textbf{Producto punto deducido en (c):}
          \[
          \vec{A}\cdot\vec{B} = A_i B_j (\hat{e}_i\cdot\hat{e}_j) = \delta_{ij} A_i B_j = \mathbf{A_i B_i}
          \]
          \item \textbf{Traza (dimensión de $\R^3$):}
          \[
          \delta_{ii} = \delta_{11} + \delta_{22} + \delta_{33} = \mathbf{3}
          \]
        \end{itemize}
      \end{exampleblock}
    \end{column}
  \end{columns}
\end{frame}

% ---------------------------------------------------------------------
\begin{frame}{Actividad 2: El Producto Cruz sin Determinantes}
  \begin{block}{Reto en clase: ¿De dónde sale el símbolo de Levi-Civita?}
    \small
    Olvida la regla del determinante $3 \times 3$. Multipliquemos $\vec{A} = A_j\hat{e}_j$ y $\vec{B} = B_k\hat{e}_k$ usando solo la propiedad distributiva:
    \[
    \vec{A} \times \vec{B} = (A_j\hat{e}_j) \times (B_k\hat{e}_k) = A_j B_k (\hat{e}_j \times \hat{e}_k).
    \]
    \begin{enumerate}
      \item[(a)] Proyecta con $\hat{e}_i$ para extraer la $i$-ésima componente de $\vec{A}\times\vec{B}$:
      \[
      (\vec{A}\times\vec{B})_i = \hat{e}_i \cdot (\vec{A}\times\vec{B}) = A_j B_k \big[\hat{e}_i \cdot (\hat{e}_j \times \hat{e}_k)\big].
      \]
      \item[(b)] Llamemos $\mathbf{\epsilon_{ijk} \equiv \hat{e}_i \cdot (\hat{e}_j \times \hat{e}_k)}$. Geométricamente, ¿qué mide el triple producto escalar de tres vectores?
      \item[(c)] Usando la interpretación de volumen de (b), responde sin hacer cuentas: ¿Por qué $\epsilon_{112} = 0$? ¿Cuánto valen $\epsilon_{123}$ y $\epsilon_{132}$?
    \end{enumerate}
  \end{block}
\end{frame}

% ---------------------------------------------------------------------
\begin{frame}{Discusión 2: $\epsilon_{ijk}$ como Volumen Orientado de la Base}
  \begin{columns}[c]
    \begin{column}{0.56\textwidth}
      \begin{block}{Construcción Geométrica de $\epsilon_{ijk}$}
        \small
        $\mathbf{\epsilon_{ijk} \equiv \hat{e}_i \cdot (\hat{e}_j \times \hat{e}_k)} = \text{Vol}(\hat{e}_i, \hat{e}_j, \hat{e}_k)$
      \end{block}
      \small
      \begin{itemize}
        \item \textbf{Si dos índices se repiten} ($\epsilon_{112}$): dos aristas paralelas $\implies$ caja aplastada de \textbf{volumen $0$} (21 casos).
        \item \textbf{Orden cíclico horario} $(1,2,3), (2,3,1), (3,1,2)$: cubo de mano derecha $\implies \mathbf{+1}$.
        \item \textbf{Orden antihorario} $(1,3,2), (2,1,3), (3,2,1)$: cubo de mano izquierda $\implies \mathbf{-1}$.
      \end{itemize}
      \begin{alertblock}{Producto Cruz Universal}
        \centering $\mathbf{(\vec{A} \times \vec{B})_i = \epsilon_{ijk} A_j B_k}$
      \end{alertblock}
    \end{column}
    \begin{column}{0.41\textwidth}
      \centering
      \begin{tikzpicture}[scale=0.82]
        \draw[line width=1.5pt, color=orangeacc] (0,0) circle (1.35cm);
        \node[circle, fill=blackmain, text=white, font=\bfseries, inner sep=3pt] at (90:1.35cm) {1};
        \node[circle, fill=blackmain, text=white, font=\bfseries, inner sep=3pt] at (330:1.35cm) {2};
        \node[circle, fill=blackmain, text=white, font=\bfseries, inner sep=3pt] at (210:1.35cm) {3};
        \draw[-{Stealth[length=2.5mm]}, line width=1.5pt, color=orangeacc] (68:1.65cm) arc (68:12:1.65cm) node[midway, right=1pt, font=\scriptsize\bfseries, color=orangeacc] {$+1$};
        \draw[-{Stealth[length=2.5mm]}, line width=1.5pt, color=orangeacc] (308:1.65cm) arc (308:252:1.65cm);
        \draw[-{Stealth[length=2.5mm]}, line width=1.5pt, color=orangeacc] (188:1.65cm) arc (188:132:1.65cm);
        \node[font=\scriptsize\bfseries, align=center] at (0,0) {Volumen\\Orientado};
      \end{tikzpicture}
    \end{column}
  \end{columns}
\end{frame}

% ---------------------------------------------------------------------
\begin{frame}{Actividad 3: Construyendo la Identidad $\epsilon-\delta$ por Casos}
  \begin{block}{Reto en clase: ¿Cuándo no es cero el producto $\epsilon_{ijk}\epsilon_{imn}$?}
    \small
    Al encadenar dos productos cruz aparece $\epsilon_{ijk}\epsilon_{imn}$. Fijemos $\mathbf{i = 1}$ y analicemos $\epsilon_{1jk}\epsilon_{1mn}$:
    \begin{enumerate}
      \item[(a)] Para que $\epsilon_{1jk} \neq 0$, los índices $(j, k)$ no pueden ser $1$ ni repetirse. ¿Cuáles son las \textbf{dos únicas parejas} posibles para $(j, k)$? ¿Y para $(m, n)$?
      \item[(b)] Compara la pareja $(j, k)$ con la pareja $(m, n)$:
      \begin{itemize}
        \item \textbf{Caso 1 (Mismo orden):} Si $j = m$ y $k = n$, ¿cuánto vale el producto $\epsilon_{1jk}\epsilon_{1mn}$?
        \item \textbf{Caso 2 (Orden invertido):} Si $j = n$ y $k = m$ (con $j \neq k$), ¿cuánto vale $\epsilon_{1jk}\epsilon_{1mn}$?
      \end{itemize}
      \item[(c)] Escribe una combinación de productos de deltas de Kronecker que valga $+1$ en el Caso 1, $-1$ en el Caso 2, y $0$ en cualquier otro caso.
    \end{enumerate}
  \end{block}
\end{frame}

% ---------------------------------------------------------------------
\begin{frame}{Discusión 3: La Identidad Maestra $\epsilon-\delta$ Desmitificada}
  \small
  \textbf{Resultados de la Actividad 3:}
  \begin{itemize}
    \item Para un $i$ dado, $(j,k)$ y $(m,n)$ deben elegir los dos números restantes de $\{1, 2, 3\} \setminus \{i\}$.
    \item Si coinciden en el mismo orden ($j=m, k=n$), ambos $\epsilon$ tienen igual signo $\implies (+1)(+1) = (-1)(-1) = \mathbf{+1}$.
    \item Si están cruzados ($j=n, k=m$), tienen signos opuestos $\implies (+1)(-1) = \mathbf{-1}$.
  \end{itemize}
  \begin{block}{Identidad Maestra Universal $\epsilon-\delta$}
    \centering $\mathbf{\epsilon_{ijk}\epsilon_{imn} = \delta_{jm}\delta_{kn} - \delta_{jn}\delta_{km}}$
  \end{block}
  \begin{exampleblock}{Lectura Geométrica}
    \begin{itemize}
      \item Primer término ($\delta_{jm}\delta_{kn}$): emparejamiento directo (primero con primero, segundo con segundo).
      \item Segundo término ($-\delta_{jn}\delta_{km}$): emparejamiento invertido (un intercambio de orientación $\implies$ signo menos).
    \end{itemize}
  \end{exampleblock}
\end{frame}

% ---------------------------------------------------------------------
\begin{frame}{Actividad 4: La Gran Demolición en 3 Renglones}
  \begin{block}{Reto final: Aplica tus herramientas}
    \small
    \begin{enumerate}
      \item[(a)] \textbf{Demuele la identidad BAC--CAB:}
      Completa los tres pasos para calcular la componente $i$ de $\vec{A} \times (\vec{B} \times \vec{C})$:
      \begin{align*}
        [\vec{A} \times (\vec{B} \times \vec{C})]_i &= \epsilon_{ijk} A_j (\vec{B}\times\vec{C})_k = \epsilon_{kij}\epsilon_{kmn} A_j B_m C_n \\
        &= (\delta_{im}\delta_{jn} - \delta_{in}\delta_{jm}) A_j B_m C_n = \dots
      \end{align*}
      \item[(b)] \textbf{Física (Fuerza Magnética de Lorentz):}
      La fuerza magnética es $F_i = q\,\epsilon_{ijk} v_j B_k$. La potencia mecánica entregada a la partícula es $P = \vec{F}\cdot\vec{v} = F_i v_i = q\,\epsilon_{ijk} v_i v_j B_k$.
      ¿Qué ocurre con la expresión $\epsilon_{ijk} v_i v_j$ si intercambias el nombre de los índices mudos $i \leftrightarrow j$?
    \end{enumerate}
  \end{block}
\end{frame}

% ---------------------------------------------------------------------
\begin{frame}{Discusión 4: Solución y Puente hacia la Semana 02}
  \begin{columns}[c]
    \begin{column}{0.52\textwidth}
      \begin{exampleblock}{(a) BAC--CAB en 3 líneas}
        \small
        \begin{align*}
          [\vec{A}\times(\vec{B}\times\vec{C})]_i &= (\delta_{im}\delta_{jn} - \delta_{in}\delta_{jm}) A_j B_m C_n \\
          &= B_i(A_j C_j) - C_i(A_j B_j) \\
          &= [\vec{B}(\vec{A}\cdot\vec{C}) - \vec{C}(\vec{A}\cdot\vec{B})]_i \quad \blacksquare
        \end{align*}
        ¡El término $i=j$ que antes sumábamos y restábamos a mano se cancela solo!
      \end{exampleblock}
    \end{column}
    \begin{column}{0.45\textwidth}
      \begin{block}{(b) Simetría vs. Antisimetría}
        \small
        Como $v_i v_j = v_j v_i$ (simétrico) y $\epsilon_{jik} = -\epsilon_{ijk}$ (antisimétrico):
        \[
        \epsilon_{ijk} v_i v_j = 0 \implies \mathbf{P = 0}
        \]
        ¡El campo magnético nunca realiza trabajo!
      \end{block}
      \vspace{0.1cm}
      \centering
      \textcolor{orangeacc}{\textbf{Siguiente parada (Semana 02):}}\\
      {\footnotesize ¿Qué pasa con $\delta_{ij}$ y $\epsilon_{ijk}$ cuando rotamos nuestro laboratorio? El grupo $SO(3)$.}
    \end{column}
  \end{columns}
\end{frame}

\end{document}
